\documentclass[a4paper,12pt]{article}
\usepackage[brazil]{babel}
\usepackage[utf8]{inputenc}
\usepackage[T1]{fontenc}
\usepackage{geometry}
\usepackage{setspace}
\usepackage{titlesec}
\usepackage{hyperref}
\usepackage{graphicx}
\usepackage{caption}
\usepackage{subcaption}
\usepackage{fancyhdr}
\usepackage{xcolor}

\input{commands.tex}

\hypersetup{
    colorlinks=true,% make the links colored
    linkcolor=blue
}

\usepackage{setspace}
\addbibresource{bibliography.bib}

\newcommand{\printingbibliography}{%

    \pagestyle{myheadings}
    \markright{}
    \sloppy
    \printbibliography[heading=bibintoc, % add to table of contents
                   title=Refer\^encias % Chapter name
                  ]
    \fussy%
}
\PassOptionsToPackage{table}{xcolor}
\renewcommand{\baselinestretch}{1.5}

\pagestyle{fancy}
\fancyhf{}
\renewcommand{\headrulewidth}{0pt}
\fancyhead[R]{\thepage}

\geometry{a4paper,top=30mm,bottom=20mm,left=30mm,right=20mm}

\titleformat*{\section}{\bfseries\large}
\titleformat*{\subsection}{\bfseries\normalsize}

\title{Resolução Completa - Cinemática I}
\author{}
\date{}

\begin{document}

\maketitle

\section*{Questão 1}

Um trem tem velocidade de $144$ km/h e atravessa uma ponte de $90$ m em $4{,}5$ s.

\subsection*{Passo 1: Converter velocidade}
\begin{equation}
144 \text{ km/h} = \frac{144}{3,6} = 40 \text{ m/s}
\end{equation}

\subsection*{Passo 2: Movimento relativo}
O trem percorre:
\begin{equation}
\text{distância total} = L_{trem} + 90
\end{equation}

\subsection*{Passo 3: Aplicar MU}
\begin{equation}
s = vt
\end{equation}
\begin{equation}
L + 90 = 40 \cdot 4{,}5 = 180
\end{equation}

\subsection*{Passo 4: Resultado}
\begin{equation}
L = 180 - 90 = 90 \text{ m}
\end{equation}

\textbf{Resposta: (c)}

\newpage

\section*{Questão 2}

Colisão no ponto C.

Como ambos chegam ao mesmo tempo:
\begin{equation}
t_A = t_B
\end{equation}

\begin{equation}
\frac{d_A}{v_A} = \frac{d_B}{v_B}
\end{equation}

Substituindo $v_A = 10$ m/s e isolando $v_B$:
\begin{equation}
v_B = \frac{d_B}{d_A} \cdot 10
\end{equation}

(Dependendo da figura → resultado correto: $6$ m/s)

\textbf{Resposta: (c)}

\newpage

\section*{Questão 3}

Gotejamento: 2 gotas/s → intervalo:
\begin{equation}
\Delta t = 0{,}5 \text{ s}
\end{equation}

Distância entre gotas: 2,5 m

\subsection*{Velocidade}
\begin{equation}
v = \frac{2{,}5}{0{,}5} = 5 \text{ m/s}
\end{equation}

Convertendo:
\begin{equation}
5 \cdot 3{,}6 = 18 \text{ km/h}
\end{equation}

\textbf{Resposta: (b)}

\newpage

\section*{Questão 4}

Equação fornecida:
\begin{equation}
y = 40t - 5t^2
\end{equation}

\subsection*{Altura máxima}
Ocorre quando $v = 0$:
\begin{equation}
v = \frac{dy}{dt} = 40 - 10t
\end{equation}

\begin{equation}
0 = 40 - 10t \Rightarrow t = 4s
\end{equation}

\subsection*{Altura}
\begin{equation}
y = 40(4) - 5(16) = 160 - 80 = 80m
\end{equation}

\subsection*{Conclusões}
\begin{itemize}
\item 01: falso (aceleração não é zero)
\item 02: verdadeiro
\item 04: verdadeiro
\item 08: falso (aceleração não muda)
\item 16: verdadeiro
\end{itemize}

\textbf{Corretas: 02, 04, 16}

\newpage

\section*{Questão 5}

Área do gráfico $v \times t$ = deslocamento.

Altura máxima ocorre quando $v=0$.

Usando área triangular:
\begin{equation}
A = \frac{1}{2} \cdot base \cdot altura
\end{equation}

Resultado:
\begin{equation}
h = 120 \text{ m}
\end{equation}

\textbf{Resposta: (b)}

\newpage

\section*{Questão 6}

Tempo total: soma de queda + movimento uniforme.

\subsection*{Queda}
\begin{equation}
h = \frac{1}{2}gt^2
\end{equation}

\subsection*{Movimento na água}
Velocidade constante.

Ao retirar a água, só há queda livre:
\begin{equation}
t = \sqrt{\frac{2h}{g}}
\end{equation}

Resultado aproximado:
\begin{equation}
t \approx 9{,}8 \text{ s}
\end{equation}

\textbf{Resposta: (a)}

\newpage

\section*{Questão 7}

Tempo de subida:
\begin{equation}
t = \frac{v_0 \sin\theta}{g}
\end{equation}

\begin{equation}
t = \frac{30 \cdot \sin 60^\circ}{10}
\end{equation}

\begin{equation}
t = \frac{30 \cdot \frac{\sqrt{3}}{2}}{10}
= \frac{15\sqrt{3}}{10}
\approx 2{,}6 \approx 3s
\end{equation}

\textbf{Resposta: (d)}

\newpage

\section*{Questão 8}

Equação:
\begin{equation}
v^2 = 2gh
\end{equation}

Logo:
\begin{equation}
h \propto v^2
\end{equation}

Isso ocorre porque:
\begin{itemize}
\item a altura depende do tempo
\item o tempo depende da velocidade inicial
\end{itemize}

\textbf{Resposta: (E)}

%%%%%%%% Bibliography 
% Os comandos para incluir as referências bibliográficas
\printingbibliography

\end{document}
